哈希表和前缀和深度题卡 按题型拆成章内大节，但每个大节都先从 LeetCode 案例切入，再进入分流、案例矩阵、案例推导和实验复盘。这样读者看到的是从具体题推到规律，而不是先读抽象方法再猜它能用在哪。

\topic{前缀和与哈希表}
这一节先看 LeetCode 1 两数之和、LeetCode 49 字母异位词分组、LeetCode 128 最长连续序列 这几道 LeetCode 题，再整理同一题型的分流和推导。阅读顺序不是先背原理，而是先看案例的暴力瓶颈，再把重复出现的观察抽象成方法。

\subsection{原理和适用条件}

以 LeetCode 1 两数之和、LeetCode 49 字母异位词分组、LeetCode 128 最长连续序列 为主线：LeetCode 1 两数之和：模型是“当前元素只需要寻找唯一补数，历史值到下标的映射就是最小状态。”，优化抓手是“先查再插入，避免同一个下标被用两次；若数组有序才考虑双指针。”；LeetCode 49 字母异位词分组：模型是“异位词共享同一个规范表示，可以是排序字符串或字符计数。”，优化抓手是“哈希表按规范键分桶，避免两两比较字符串。”；LeetCode 128 最长连续序列：模型是“连续段只需要从没有前驱的数开始扩展。”，优化抓手是“哈希集合判断 x 减一 是否存在，跳过非起点避免重复扫描。”。这些案例共同推出的规律是：先用双层枚举区间，再把区间条件改写成两个前缀状态的关系。哈希表保存历史前缀的次数、最早位置或存在性，具体保存什么由题目目标决定。 方法名只是复盘后的标签，真正要掌握的是从题面约束、暴力失败、状态设计到边界反例的推导链。

\begin{principlebox}{图示：从 LeetCode 案例到 前缀和与哈希表推导链}
\begin{tabular}{@{}p{0.18\linewidth}p{0.74\linewidth}@{}}
暴力基线 & 枚举候选、完整模拟或逐个检查，先保证覆盖全部可能答案。\\
瓶颈定位 & 瓶颈在于当前右端需要寻找很多历史左端。若这些历史状态能被同一个键表示，就不该逐个扫描，而应把历史状态压进哈希表。\\
结构选择 & 前缀和把区间问题改写成两个前缀状态的差；哈希表把寻找匹配前缀改成查表。频次表、最早位置表、最近位置表对应不同目标：计数、最长区间、最短区间不能混用。\\
不变量 & 哈希表的不变量通常是当前下标之前的历史前缀信息。先查询还是先更新不是风格问题，而是决定是否允许空区间和是否重复使用当前前缀。\\
代码落地 & 实现时用 \code{find} 判断是否存在，用 \code{operator[]} 只在确认要插入或累加时使用。总和可能溢出时使用 \code{long long}；模运算要处理负数余数归一化。\\
\end{tabular}
\end{principlebox}

\subsection{LeetCode 案例分流}

\begin{principlebox}{从 LeetCode 案例分流：前缀和与哈希表}
\textbf{补数和历史值。}
代表案例：LeetCode 1 两数之和、LeetCode 217 存在重复元素。先看这些题的题面条件：当前元素只需要寻找历史中是否存在某个配对值。由此推出的处理方式是：扫描时先查历史表，再插入当前值。风险点：先查后插能避免同一位置被重复使用。

\textbf{前缀和计数。}
代表案例：LeetCode 560 和为 K 的子数组、LeetCode 1248 统计优美子数组。先看这些题的题面条件：区间条件可以写成当前前缀和历史前缀的差。由此推出的处理方式是：哈希表保存前缀值出现次数，当前前缀查询目标历史值。风险点：计数题保存频次，最长题保存最早位置，不能混用。

\textbf{规范 key 分组。}
代表案例：LeetCode 49 字母异位词分组、LeetCode 128 最长连续序列。先看这些题的题面条件：多个对象在某种规范表示下等价。由此推出的处理方式是：把排序结果、计数向量或状态编码作为 key 分桶。风险点：key 必须无歧义，字符集和分隔符要明确。

\textbf{函数图和重复状态。}
代表案例：LeetCode 202 快乐数、LeetCode 287 寻找重复数。先看这些题的题面条件：状态经过确定性变换后可能进入循环。由此推出的处理方式是：用集合检测重复状态，或转成快慢指针。风险点：必须确认状态空间有限，且每一步变换确定。

\end{principlebox}

\subsection{案例矩阵}

本节案例覆盖 LeetCode 1 两数之和、LeetCode 49 字母异位词分组、LeetCode 128 最长连续序列、LeetCode 202 快乐数、LeetCode 217 存在重复元素、LeetCode 242 有效的字母异位词、LeetCode 303 区域和检索 - 数组不可变、LeetCode 325 和等于 k 的最长子数组长度、LeetCode 454 四数相加 II、LeetCode 523 连续的子数组和、LeetCode 525 连续数组、LeetCode 560 和为 K 的子数组、LeetCode 974 和可被 K 整除的子数组、LeetCode 1074 元素和为目标值的子矩阵数量、LeetCode 1248 统计优美子数组。阅读时不要按题号背答案，而要先判断题面落在哪个分流场景：查询目标是什么，输入约束给了什么单调性或等价关系，暴力慢在哪里，优化结构保存了什么状态。

\begin{principlebox}{案例矩阵}
\textbf{LeetCode 1 两数之和。}
当前元素只需要寻找唯一补数，历史值到下标的映射就是最小状态。单点风险：重复数字、目标为偶数且补数等于当前值、没有答案时返回空结果。

\textbf{LeetCode 49 字母异位词分组。}
异位词共享同一个规范表示，可以是排序字符串或字符计数。单点风险：空字符串、重复单词、字符集不止小写字母、计数键必须无歧义。

\textbf{LeetCode 128 最长连续序列。}
连续段只需要从没有前驱的数开始扩展。单点风险：重复元素、负数、空数组、整数边界。

\textbf{LeetCode 202 快乐数。}
反复变换数字会进入一或循环，状态空间最终有限。单点风险：输入为一、进入循环、位平方和函数、负数不在题意内。

\textbf{LeetCode 217 存在重复元素。}
判断是否出现过是哈希集合最直接的职责。单点风险：空数组、全唯一、重复在末尾、值域很小。

\textbf{LeetCode 242 有效的字母异位词。}
两个字符串是异位词等价于每个字符出现次数相同。单点风险：长度不同、空串、Unicode 字符、计数减到负数。

\textbf{LeetCode 303 区域和检索 - 数组不可变。}
数组不变且要多次区间求和，单次扫描区间会重复读取同一段。单点风险：left 为零、负数、单元素区间、多次查询、总和溢出。

\textbf{LeetCode 325 和等于 k 的最长子数组长度。}
最长区间要求在同一前缀差关系下尽量拉大左右端距离。单点风险：从下标零开始、负数、前缀重复时不能覆盖最早位置、答案不存在。

\textbf{LeetCode 454 四数相加 II。}
四个数组的和为零可以拆成两个二数组和互为相反数。单点风险：重复值、负数、答案数量很大、两数和范围。

\textbf{LeetCode 523 连续的子数组和。}
长度至少二的子数组和为 k 的倍数，等价于两个前缀和同余。单点风险：k 为零、负数余数、长度限制、从开头开始的区间。

\textbf{LeetCode 525 连续数组。}
把零看成负一后，零和一数量相等等价于前缀和相同。单点风险：全零、全一、从下标零开始、不要覆盖最早位置。

\textbf{LeetCode 560 和为 K 的子数组。}
当前前缀和减去目标值，就是需要匹配的历史前缀。单点风险：负数、目标为零、从下标零开始的区间、重复前缀。

\textbf{LeetCode 974 和可被 K 整除的子数组。}
区间和能被 K 整除等价于两个前缀和余数相同。单点风险：负数余数归一化、K 为一、从开头开始、重复余数很多。

\textbf{LeetCode 1074 元素和为目标值的子矩阵数量。}
二维矩阵子区域和可以固定上下边界后转成一维子数组和。单点风险：单行、单列、负数、目标为零、矩阵较大。

\textbf{LeetCode 1248 统计优美子数组。}
恰好 K 个奇数可转成前缀奇数计数差为 K。单点风险：K 为零、全偶数、全奇数、从开头开始的区间。

\end{principlebox}

\subsection{共性落点}

\begin{principlebox}{本组共性落点}
\textbf{慢版本。} 暴力做法枚举所有区间，直接求区间和或重新比较频次。即使用前缀和把求和降到常数，仍可能保留枚举左右端点的平方复杂度。
\textbf{瓶颈。} 瓶颈在于当前右端需要寻找很多历史左端。若这些历史状态能被同一个键表示，就不该逐个扫描，而应把历史状态压进哈希表。
\textbf{状态字段。} 当前状态或规范键、需要查询的历史 key、哈希表 value 语义、答案累计值或分桶结果；空前缀、空字符串、重复状态或初始状态要单独说明。
\textbf{不变量。} 哈希表的不变量通常是当前下标之前的历史前缀信息。先查询还是先更新不是风格问题，而是决定是否允许空区间和是否重复使用当前前缀。
\textbf{实现检查。} 先查询还是先插入要由题意决定；key 构造必须无歧义；哈希表 value 的默认插入副作用要可控。
\textbf{C++ 落地。} 实现时用 \code{find} 判断是否存在，用 \code{operator[]} 只在确认要插入或累加时使用。总和可能溢出时使用 \code{long long}；模运算要处理负数余数归一化。
\textbf{证明抓手。} 证明从状态等价开始：任意合法答案都对应当前状态和某个历史 key、规范键或访问状态，查到的 key 也能反推出合法答案或合法分组。
\end{principlebox}

\subsection{案例推导}

\noindent\textbf{LeetCode 1 两数之和。}

\textbf{题目说明。} LeetCode 1 两数之和 的题面入口要先还原成具体任务：当前元素只需要寻找唯一补数，历史值到下标的映射就是最小状态。

\textbf{样例入口。} 拿数组 [2, 7]、target=9 手算：站在 2 时历史为空，不能配对；站在 7 时只需要问历史里有没有 2。再试 [3, 3]、target=6：如果先把当前 3 插入再查，就会把同一个下标用两次。

\textbf{暴力基线。} 暴力两层循环枚举 i 和 j，检查 nums[i]+nums[j] 是否等于 target，找到后返回两个下标。

\textbf{暴力过程。} 以 nums=[2, 7, 11, 15]、target=9 为例，暴力先试 2+7 命中；换成 [3, 2, 4]、target=6 时，会试 3+2 失败，再试 2+4 命中。

\textbf{重复工作。} 题面模型：当前元素只需要寻找唯一补数，历史值到下标的映射就是最小状态。暴力版的慢点要落到本题：浪费在每个当前位置都回头扫描历史。站在 nums[i] 时，真正需要的不是所有历史值，而是历史里有没有 target-nums[i]。后面的优化只消掉这类重复工作，不能改变答案集合。

\textbf{优化条件。} 本题的关键观察是：先查再插入，避免同一个下标被用两次；若数组有序才考虑双指针。这句话必须能翻译成可维护状态；如果题目条件改变后观察不再成立，就不能继续套原来的写法。

\textbf{相邻状态。} 规律不是“用哈希表”四个字，而是每个当前位置只和它左侧历史配对，代码必须先查补数再插入当前值。把这个特例继续拆开看：每个合法答案都要还原成当前状态和某个历史状态的配对；哈希表的 key 负责定位历史状态，value 负责表达次数、最早位置或存在性。先查还是先写，必须由是否允许当前前缀和自己配对来决定。

\textbf{状态复用。} 哈希表保存已经扫描过的值到下标的映射。每到一个新数，先查补数是否存在；不存在再把当前数放入历史。手算时只改被新输入影响的那一小部分，而不是回头重算完整候选。

\textbf{指针和字段。} 状态字段要能说清：i 是当前下标，need=target-nums[i] 是补数，seen[value] 是该值在左侧历史中出现的下标。手算时按这个协议跟踪：nums=[3, 3]、target=6 时，第一个 3 先查不到补数，再插入；第二个 3 查到历史 3，返回两个不同下标。

\textbf{代码落点。} 关键顺序是先 find(need)，再 seen[nums[i]]=i。若先插入当前值，target=2*nums[i] 时可能把同一位置用两次。关键代码行必须能逐句翻译成“旧状态怎样变成新状态”，否则就只是模板。

\textbf{正确性。} 每个答案对都有一个较右下标；扫描到较右下标时，较左下标已经在 seen 中，所以一定能被查到。

\textbf{边界实验。} 先用重复数字、目标为偶数且补数等于当前值、没有答案时返回空结果做反例检查。实验时把重复数字、目标为偶数且补数等于当前值、没有答案时返回空结果加入对拍集合，用平方暴力和优化版比较。失败时先看初始历史状态、查询更新顺序和哈希表 value 类型。

\textbf{复杂度变化。} 暴力 O(\ensuremath{n^{2}})，哈希表扫描 O(n) 均摊时间、O(n) 空间。

\textbf{迁移追问。} 如果题目改成“若改成返回所有不重复数对，要先排序或用频次表处理重复”，先判断原状态是否仍然覆盖新答案集合，再决定是微调边界、改 value 语义，还是换算法。

\noindent\textbf{LeetCode 49 字母异位词分组。}

\textbf{题目说明。} LeetCode 49 字母异位词分组 的题面入口要先还原成具体任务：异位词共享同一个规范表示，可以是排序字符串或字符计数。

\textbf{样例入口。} 拿 eat、tea、tan 手算，eat 和 tea 排序后都是 aet，或计数向量都是 a1e1t1；tan 的键不同。两两比较会重复比较很多字符串，哈希分桶只比较规范键。

\textbf{暴力基线。} 暴力两两比较字符串是否互为异位词，或者对每个新字符串扫描已有分组找能放入的组。

\textbf{暴力过程。} 以 eat、tea、tan 为例，eat 和 tea 排序后都是 aet，应放一组；tan 排序后是 ant，属于另一组。

\textbf{重复工作。} 题面模型：异位词共享同一个规范表示，可以是排序字符串或字符计数。暴力版的慢点要落到本题：浪费在重复比较字符串。异位词的本质是字符频次相同，只要把每个字符串转成规范 key，就能一次定位分组。后面的优化只消掉这类重复工作，不能改变答案集合。

\textbf{优化条件。} 本题的关键观察是：哈希表按规范键分桶，避免两两比较字符串。这句话必须能翻译成可维护状态；如果题目条件改变后观察不再成立，就不能继续套原来的写法。

\textbf{相邻状态。} 规律是异位词共享规范表示，字符集固定时计数键更直接；空字符串和重复单词也必须落到同一个键语义里。把这个特例继续拆开看：每个合法答案都要还原成当前状态和某个历史状态的配对；哈希表的 key 负责定位历史状态，value 负责表达次数、最早位置或存在性。先查还是先写，必须由是否允许当前前缀和自己配对来决定。

\textbf{状态复用。} 哈希表 group[key] 保存同一个规范 key 下的所有原字符串。key 可以是排序后的字符串，也可以是 26 维计数向量编码。手算时只改被新输入影响的那一小部分，而不是回头重算完整候选。

\textbf{指针和字段。} 状态字段要能说清：word 是当前字符串，key 是它的规范表示，groups[key] 是该异位词类的结果数组。手算时按这个协议跟踪：扫描 eat 时创建 key=aet 的组；扫描 tea 时得到同一个 key，追加进同组；扫描 tan 时进入 key=ant 的新组。

\textbf{代码落点。} 排序 key 写法简单但每个词 O(L log L)；计数 key 是 O(L)，但编码时必须加分隔符，避免 1 和 11 这类计数串歧义。关键代码行必须能逐句翻译成“旧状态怎样变成新状态”，否则就只是模板。

\textbf{正确性。} 两个字符串互为异位词当且仅当字符频次完全相同；规范 key 完整表达频次，所以同 key 等价、不同 key 不等价。

\textbf{边界实验。} 先用空字符串、重复单词、字符集不止小写字母、计数键必须无歧义做反例检查。实验时把空字符串、重复单词、字符集不止小写字母、计数键必须无歧义加入对拍集合，用平方暴力和优化版比较。失败时先看初始历史状态、查询更新顺序和哈希表 value 类型。

\textbf{复杂度变化。} 两两比较最坏 O(\ensuremath{n^{2}}L)；排序 key 为 O(nL log L)，计数 key 为 O(nL)，空间 O(nL)。

\textbf{迁移追问。} 如果题目改成“若字符串很长且字符集固定，计数键通常比排序键更稳定”，先判断原状态是否仍然覆盖新答案集合，再决定是微调边界、改 value 语义，还是换算法。

\noindent\textbf{LeetCode 128 最长连续序列。}

\textbf{题目说明。} LeetCode 128 最长连续序列 的题面入口要先还原成具体任务：连续段只需要从没有前驱的数开始扩展。

\textbf{样例入口。} 拿 [100, 4, 200, 1, 3, 2] 手算，4 不是起点，因为 3 存在；1 没有前驱，从 1 扩到 4 得到长度 4。

\textbf{暴力基线。} 暴力对每个数 x，向上查找 x+1、x+2 是否存在，计算以 x 开头的连续长度。

\textbf{暴力过程。} 以 [100, 4, 200, 1, 3, 2] 为例，若从 2、3、4 都向后扩，会反复扫描同一段 1-4。

\textbf{重复工作。} 题面模型：连续段只需要从没有前驱的数开始扩展。暴力版的慢点要落到本题：浪费在每个连续段被多个内部元素重复扩展。一个连续段只应该从没有前驱 x-1 的起点开始扩。后面的优化只消掉这类重复工作，不能改变答案集合。

\textbf{优化条件。} 本题的关键观察是：哈希集合判断 x 减一 是否存在，跳过非起点避免重复扫描。这句话必须能翻译成可维护状态；如果题目条件改变后观察不再成立，就不能继续套原来的写法。

\textbf{相邻状态。} 规律是只从没有 x-1 的数开始扩展，避免每个连续段被重复扫描；重复元素先由集合去掉。把这个特例继续拆开看：每个合法答案都要还原成当前状态和某个历史状态的配对；哈希表的 key 负责定位历史状态，value 负责表达次数、最早位置或存在性。先查还是先写，必须由是否允许当前前缀和自己配对来决定。

\textbf{状态复用。} 哈希集合保存所有数字。遍历每个 x，只有当 x-1 不存在时，才从 x 开始不断找 x+1。手算时只改被新输入影响的那一小部分，而不是回头重算完整候选。

\textbf{指针和字段。} 状态字段要能说清：num 是候选起点，current 是扩展中的数，length 是当前连续段长度，best 是最大长度。手算时按这个协议跟踪：4 有前驱 3，不做起点；1 没有前驱，从 1 扩到 2、3、4，长度 4。

\textbf{代码落点。} 关键是 if (!set.contains(x-1)) 才进入 while。重复元素先被集合去掉，不影响长度。关键代码行必须能逐句翻译成“旧状态怎样变成新状态”，否则就只是模板。

\textbf{正确性。} 每个连续段有且只有一个没有前驱的最小元素；算法只从这个元素扩展，覆盖每段且不重复扩展内部点。

\textbf{边界实验。} 先用重复元素、负数、空数组、整数边界做反例检查。实验时把重复元素、负数、空数组、整数边界加入对拍集合，用平方暴力和优化版比较。失败时先看初始历史状态、查询更新顺序和哈希表 value 类型。

\textbf{复杂度变化。} 排序法 O(n log n)；哈希起点法均摊 O(n) 时间、O(n) 空间。

\textbf{迁移追问。} 如果题目改成“若数据流动态插入，可以用并查集或区间合并维护长度”，先判断原状态是否仍然覆盖新答案集合，再决定是微调边界、改 value 语义，还是换算法。

\noindent\textbf{LeetCode 202 快乐数。}

\textbf{题目说明。} LeetCode 202 快乐数 的题面入口要先还原成具体任务：反复变换数字会进入一或循环，状态空间最终有限。

\textbf{样例入口。} 拿 19 手算，1 的平方加 9 的平方得到 82，82 再变成 68、100、1，所以快乐；拿 2 会进入循环而不到 1。

\textbf{暴力基线。} 慢版本按题意两两比较、枚举区间、扫描历史或持续迭代并查重。把检查条件写出来后，通常能看到当前状态只缺一个历史状态、规范键或已经访问过的中间状态；这正是从平方枚举或重复迭代走向哈希表的入口。放到本题就是：先按“反复变换数字会进入一或循环，状态空间最终有限”完整试慢版本；样例里已经能看到“规律是用集合记录出现过的中间值，或用快慢指针检测平方和序列是否成环”，所以优化前必须先能说清慢版本枚举了哪些候选。

\textbf{暴力过程。} 拿 19 手算，1 的平方加 9 的平方得到 82，82 再变成 68、100、1，所以快乐；拿 2 会进入循环而不到 1。

\textbf{重复工作。} 题面模型：反复变换数字会进入一或循环，状态空间最终有限。暴力版的慢点要落到本题：慢版本会为“反复变换数字会进入一或循环，状态空间最终有限”反复寻找历史状态；而样例规律已经指出“规律是用集合记录出现过的中间值，或用快慢指针检测平方和序列是否成环”。这些补数、前缀状态、规范键、半边和或访问过的中间状态可以用一个 key 归类，不必每次重新扫描或重复比较。后面的优化只消掉这类重复工作，不能改变答案集合。

\textbf{优化条件。} 本题的关键观察是：用集合检测重复状态，或用快慢指针看成函数图。这句话必须能翻译成可维护状态；如果题目条件改变后观察不再成立，就不能继续套原来的写法。

\textbf{相邻状态。} 规律是用集合记录出现过的中间值，或用快慢指针检测平方和序列是否成环。把这个特例继续拆开看：每个合法答案都要还原成当前状态和某个历史状态的配对；哈希表的 key 负责定位历史状态，value 负责表达次数、最早位置或存在性。先查还是先写，必须由是否允许当前前缀和自己配对来决定。

\textbf{状态复用。} 把回头扫描、两两比较或重复状态检测改成查表、分桶或访问记录；key 是当前题目需要的补数、前缀状态、规范表示、半边和或中间状态，value 根据目标选择次数、下标、列表或存在性。本题落点是：规律是用集合记录出现过的中间值，或用快慢指针检测平方和序列是否成环。手算时只改被新输入影响的那一小部分，而不是回头重算完整候选。

\textbf{指针和字段。} 状态字段要能说清：当前状态或规范键、需要查询的历史 key、哈希表 value 语义、答案累计值或分桶结果；空前缀、空字符串、重复状态或初始状态要单独说明。手算时按这个协议跟踪：拿 19 手算，1 的平方加 9 的平方得到 82，82 再变成 68、100、1，所以快乐；拿 2 会进入循环而不到 1。然后按协议复查：手算时列出当前输入、当前状态或规范键、需要查询的历史 key、查到的 value、更新后的哈希表。这样能看出先查后插、先分桶后输出，还是先检测重复状态。

\textbf{代码落点。} 先查询还是先插入要由题意决定；key 构造必须无歧义；哈希表 value 的默认插入副作用要可控。关键代码行必须能逐句翻译成“旧状态怎样变成新状态”，否则就只是模板。

\textbf{正确性。} 证明从状态等价开始：任意合法答案都对应当前状态和某个历史 key、规范键或访问状态，查到的 key 也能反推出合法答案或合法分组。

\textbf{边界实验。} 先用输入为一、进入循环、位平方和函数、负数不在题意内做反例检查。实验时把输入为一、进入循环、位平方和函数、负数不在题意内加入对拍集合，用平方暴力和优化版比较。失败时先看初始历史状态、查询更新顺序和哈希表 value 类型。

\textbf{复杂度变化。} 暴力两两比较、枚举区间、扫描历史或重复迭代常见为平方级或更多；把历史状态、规范键或访问状态放进哈希表后，每步只做常数次查询和更新，通常是线性时间、线性空间。

\textbf{迁移追问。} 如果题目改成“快慢指针版本能训练把数值迭代看成链表”，先判断原状态是否仍然覆盖新答案集合，再决定是微调边界、改 value 语义，还是换算法。

\noindent\textbf{LeetCode 217 存在重复元素。}

\textbf{题目说明。} LeetCode 217 存在重复元素 的题面入口要先还原成具体任务：判断是否出现过是哈希集合最直接的职责。

\textbf{样例入口。} 拿 [1, 2, 3, 1] 手算，扫描到最后一个 1 时，哈希集合里已经有 1，立即返回重复。

\textbf{暴力基线。} 慢版本按题意两两比较、枚举区间、扫描历史或持续迭代并查重。把检查条件写出来后，通常能看到当前状态只缺一个历史状态、规范键或已经访问过的中间状态；这正是从平方枚举或重复迭代走向哈希表的入口。放到本题就是：先按“判断是否出现过是哈希集合最直接的职责”完整试慢版本；样例里已经能看到“规律是集合保存已经出现过的历史值，插入前查询可以在第一次重复处停止；空数组和无重复数组返回假”，所以优化前必须先能说清慢版本枚举了哪些候选。

\textbf{暴力过程。} 拿 [1, 2, 3, 1] 手算，扫描到最后一个 1 时，哈希集合里已经有 1，立即返回重复。

\textbf{重复工作。} 题面模型：判断是否出现过是哈希集合最直接的职责。暴力版的慢点要落到本题：慢版本会为“判断是否出现过是哈希集合最直接的职责”反复寻找历史状态；而样例规律已经指出“规律是集合保存已经出现过的历史值，插入前查询可以在第一次重复处停止；空数组和无重复数组返回假”。这些补数、前缀状态、规范键、半边和或访问过的中间状态可以用一个 key 归类，不必每次重新扫描或重复比较。后面的优化只消掉这类重复工作，不能改变答案集合。

\textbf{优化条件。} 本题的关键观察是：扫描时若插入前已经存在则立即返回真。这句话必须能翻译成可维护状态；如果题目条件改变后观察不再成立，就不能继续套原来的写法。

\textbf{相邻状态。} 规律是集合保存已经出现过的历史值，插入前查询可以在第一次重复处停止；空数组和无重复数组返回假。把这个特例继续拆开看：每个合法答案都要还原成当前状态和某个历史状态的配对；哈希表的 key 负责定位历史状态，value 负责表达次数、最早位置或存在性。先查还是先写，必须由是否允许当前前缀和自己配对来决定。

\textbf{状态复用。} 把回头扫描、两两比较或重复状态检测改成查表、分桶或访问记录；key 是当前题目需要的补数、前缀状态、规范表示、半边和或中间状态，value 根据目标选择次数、下标、列表或存在性。本题落点是：规律是集合保存已经出现过的历史值，插入前查询可以在第一次重复处停止；空数组和无重复数组返回假。手算时只改被新输入影响的那一小部分，而不是回头重算完整候选。

\textbf{指针和字段。} 状态字段要能说清：当前状态或规范键、需要查询的历史 key、哈希表 value 语义、答案累计值或分桶结果；空前缀、空字符串、重复状态或初始状态要单独说明。手算时按这个协议跟踪：拿 [1, 2, 3, 1] 手算，扫描到最后一个 1 时，哈希集合里已经有 1，立即返回重复。然后按协议复查：手算时列出当前输入、当前状态或规范键、需要查询的历史 key、查到的 value、更新后的哈希表。这样能看出先查后插、先分桶后输出，还是先检测重复状态。

\textbf{代码落点。} 先查询还是先插入要由题意决定；key 构造必须无歧义；哈希表 value 的默认插入副作用要可控。关键代码行必须能逐句翻译成“旧状态怎样变成新状态”，否则就只是模板。

\textbf{正确性。} 证明从状态等价开始：任意合法答案都对应当前状态和某个历史 key、规范键或访问状态，查到的 key 也能反推出合法答案或合法分组。

\textbf{边界实验。} 先用空数组、全唯一、重复在末尾、值域很小做反例检查。实验时把空数组、全唯一、重复在末尾、值域很小加入对拍集合，用平方暴力和优化版比较。失败时先看初始历史状态、查询更新顺序和哈希表 value 类型。

\textbf{复杂度变化。} 暴力两两比较、枚举区间、扫描历史或重复迭代常见为平方级或更多；把历史状态、规范键或访问状态放进哈希表后，每步只做常数次查询和更新，通常是线性时间、线性空间。

\textbf{迁移追问。} 如果题目改成“若允许修改输入，排序后看相邻可以降低额外空间”，先判断原状态是否仍然覆盖新答案集合，再决定是微调边界、改 value 语义，还是换算法。

\noindent\textbf{LeetCode 242 有效的字母异位词。}

\textbf{题目说明。} LeetCode 242 有效的字母异位词 的题面入口要先还原成具体任务：两个字符串是异位词等价于每个字符出现次数相同。

\textbf{样例入口。} 拿 anagram 和 nagaram 手算，两个字符串每个字符计数都相同；拿 rat 和 car，r 和 c 的计数不同。

\textbf{暴力基线。} 慢版本按题意两两比较、枚举区间、扫描历史或持续迭代并查重。把检查条件写出来后，通常能看到当前状态只缺一个历史状态、规范键或已经访问过的中间状态；这正是从平方枚举或重复迭代走向哈希表的入口。放到本题就是：先按“两个字符串是异位词等价于每个字符出现次数相同”完整试慢版本；样例里已经能看到“规律是异位词比较的是字符频次，不是字符顺序；字符集固定可用数组，泛化字符集用哈希表”，所以优化前必须先能说清慢版本枚举了哪些候选。

\textbf{暴力过程。} 拿 anagram 和 nagaram 手算，两个字符串每个字符计数都相同；拿 rat 和 car，r 和 c 的计数不同。

\textbf{重复工作。} 题面模型：两个字符串是异位词等价于每个字符出现次数相同。暴力版的慢点要落到本题：慢版本会为“两个字符串是异位词等价于每个字符出现次数相同”反复寻找历史状态；而样例规律已经指出“规律是异位词比较的是字符频次，不是字符顺序；字符集固定可用数组，泛化字符集用哈希表”。这些补数、前缀状态、规范键、半边和或访问过的中间状态可以用一个 key 归类，不必每次重新扫描或重复比较。后面的优化只消掉这类重复工作，不能改变答案集合。

\textbf{优化条件。} 本题的关键观察是：固定小写字母集时用定长计数数组，字符集未知时用哈希表。这句话必须能翻译成可维护状态；如果题目条件改变后观察不再成立，就不能继续套原来的写法。

\textbf{相邻状态。} 规律是异位词比较的是字符频次，不是字符顺序；字符集固定可用数组，泛化字符集用哈希表。把这个特例继续拆开看：每个合法答案都要还原成当前状态和某个历史状态的配对；哈希表的 key 负责定位历史状态，value 负责表达次数、最早位置或存在性。先查还是先写，必须由是否允许当前前缀和自己配对来决定。

\textbf{状态复用。} 把回头扫描、两两比较或重复状态检测改成查表、分桶或访问记录；key 是当前题目需要的补数、前缀状态、规范表示、半边和或中间状态，value 根据目标选择次数、下标、列表或存在性。本题落点是：规律是异位词比较的是字符频次，不是字符顺序；字符集固定可用数组，泛化字符集用哈希表。手算时只改被新输入影响的那一小部分，而不是回头重算完整候选。

\textbf{指针和字段。} 状态字段要能说清：当前状态或规范键、需要查询的历史 key、哈希表 value 语义、答案累计值或分桶结果；空前缀、空字符串、重复状态或初始状态要单独说明。手算时按这个协议跟踪：拿 anagram 和 nagaram 手算，两个字符串每个字符计数都相同；拿 rat 和 car，r 和 c 的计数不同。然后按协议复查：手算时列出当前输入、当前状态或规范键、需要查询的历史 key、查到的 value、更新后的哈希表。这样能看出先查后插、先分桶后输出，还是先检测重复状态。

\textbf{代码落点。} 先查询还是先插入要由题意决定；key 构造必须无歧义；哈希表 value 的默认插入副作用要可控。关键代码行必须能逐句翻译成“旧状态怎样变成新状态”，否则就只是模板。

\textbf{正确性。} 证明从状态等价开始：任意合法答案都对应当前状态和某个历史 key、规范键或访问状态，查到的 key 也能反推出合法答案或合法分组。

\textbf{边界实验。} 先用长度不同、空串、Unicode 字符、计数减到负数做反例检查。实验时把长度不同、空串、Unicode 字符、计数减到负数加入对拍集合，用平方暴力和优化版比较。失败时先看初始历史状态、查询更新顺序和哈希表 value 类型。

\textbf{复杂度变化。} 暴力两两比较、枚举区间、扫描历史或重复迭代常见为平方级或更多；把历史状态、规范键或访问状态放进哈希表后，每步只做常数次查询和更新，通常是线性时间、线性空间。

\textbf{迁移追问。} 如果题目改成“异位词分组把这个二元比较扩展成规范 key 分桶”，先判断原状态是否仍然覆盖新答案集合，再决定是微调边界、改 value 语义，还是换算法。

\noindent\textbf{LeetCode 303 区域和检索 - 数组不可变。}

\textbf{题目说明。} LeetCode 303 区域和检索 - 数组不可变 的题面入口要先还原成具体任务：数组不变且要多次区间求和，单次扫描区间会重复读取同一段。

\textbf{样例入口。} 拿 nums=[-2, 0, 3, -5, 2, -1] 手算，sumRange(0, 2)=1，sumRange(2, 5)=-1。若每次查询都扫描区间会重复读同一段；预处理 prefix 后，闭区间和等于 prefix[right+1]-prefix[left]。

\textbf{暴力基线。} 慢版本按题意两两比较、枚举区间、扫描历史或持续迭代并查重。把检查条件写出来后，通常能看到当前状态只缺一个历史状态、规范键或已经访问过的中间状态；这正是从平方枚举或重复迭代走向哈希表的入口。放到本题就是：先按“数组不变且要多次区间求和，单次扫描区间会重复读取同一段”完整试慢版本；样例里已经能看到“规律是数组不可变时把所有前缀和提前缓存”，所以优化前必须先能说清慢版本枚举了哪些候选。

\textbf{暴力过程。} 拿 nums=[-2, 0, 3, -5, 2, -1] 手算，sumRange(0, 2)=1，sumRange(2, 5)=-1。若每次查询都扫描区间会重复读同一段；预处理 prefix 后，闭区间和等于 prefix[right+1]-prefix[left]。

\textbf{重复工作。} 题面模型：数组不变且要多次区间求和，单次扫描区间会重复读取同一段。暴力版的慢点要落到本题：慢版本会为“数组不变且要多次区间求和，单次扫描区间会重复读取同一段”反复寻找历史状态；而样例规律已经指出“规律是数组不可变时把所有前缀和提前缓存”。这些补数、前缀状态、规范键、半边和或访问过的中间状态可以用一个 key 归类，不必每次重新扫描或重复比较。后面的优化只消掉这类重复工作，不能改变答案集合。

\textbf{优化条件。} 本题的关键观察是：预处理 prefix[i+1] 为前 i+1 个数的和，sumRange(left, right) 用两个前缀相减。这句话必须能翻译成可维护状态；如果题目条件改变后观察不再成立，就不能继续套原来的写法。

\textbf{相邻状态。} 规律是数组不可变时把所有前缀和提前缓存。把这个特例继续拆开看：每个合法答案都要还原成当前状态和某个历史状态的配对；哈希表的 key 负责定位历史状态，value 负责表达次数、最早位置或存在性。先查还是先写，必须由是否允许当前前缀和自己配对来决定。

\textbf{状态复用。} 把回头扫描、两两比较或重复状态检测改成查表、分桶或访问记录；key 是当前题目需要的补数、前缀状态、规范表示、半边和或中间状态，value 根据目标选择次数、下标、列表或存在性。本题落点是：规律是数组不可变时把所有前缀和提前缓存。手算时只改被新输入影响的那一小部分，而不是回头重算完整候选。

\textbf{指针和字段。} 状态字段要能说清：当前状态或规范键、需要查询的历史 key、哈希表 value 语义、答案累计值或分桶结果；空前缀、空字符串、重复状态或初始状态要单独说明。手算时按这个协议跟踪：拿 nums=[-2, 0, 3, -5, 2, -1] 手算，sumRange(0, 2)=1，sumRange(2, 5)=-1。若每次查询都扫描区间会重复读同一段；预处理 prefix 后，闭区间和等于 prefix[right+1]-prefix[left]。然后按协议复查：手算时列出当前输入、当前状态或规范键、需要查询的历史 key、查到的 value、更新后的哈希表。这样能看出先查后插、先分桶后输出，还是先检测重复状态。

\textbf{代码落点。} 先查询还是先插入要由题意决定；key 构造必须无歧义；哈希表 value 的默认插入副作用要可控。关键代码行必须能逐句翻译成“旧状态怎样变成新状态”，否则就只是模板。

\textbf{正确性。} 证明从状态等价开始：任意合法答案都对应当前状态和某个历史 key、规范键或访问状态，查到的 key 也能反推出合法答案或合法分组。

\textbf{边界实验。} 先用left 为零、负数、单元素区间、多次查询、总和溢出做反例检查。实验时把left 为零、负数、单元素区间、多次查询、总和溢出加入对拍集合，用平方暴力和优化版比较。失败时先看初始历史状态、查询更新顺序和哈希表 value 类型。

\textbf{复杂度变化。} 暴力两两比较、枚举区间、扫描历史或重复迭代常见为平方级或更多；把历史状态、规范键或访问状态放进哈希表后，每步只做常数次查询和更新，通常是线性时间、线性空间。

\textbf{迁移追问。} 如果题目改成“若数组会更新，普通前缀和失效，需要树状数组或线段树”，先判断原状态是否仍然覆盖新答案集合，再决定是微调边界、改 value 语义，还是换算法。

\noindent\textbf{LeetCode 325 和等于 k 的最长子数组长度。}

\textbf{题目说明。} LeetCode 325 和等于 k 的最长子数组长度 的题面入口要先还原成具体任务：最长区间要求在同一前缀差关系下尽量拉大左右端距离。

\textbf{样例入口。} 拿 [1, -1, 5, -2, 3]、k=3 手算，到前缀和 3 时需要历史前缀 0；到前缀和 4 时需要历史前缀 1，最早的 1 能给出更长区间。

\textbf{暴力基线。} 慢版本按题意两两比较、枚举区间、扫描历史或持续迭代并查重。把检查条件写出来后，通常能看到当前状态只缺一个历史状态、规范键或已经访问过的中间状态；这正是从平方枚举或重复迭代走向哈希表的入口。放到本题就是：先按“最长区间要求在同一前缀差关系下尽量拉大左右端距离”完整试慢版本；样例里已经能看到“规律是最长区间要保存每个前缀和的最早下标，而不是最新下标”，所以优化前必须先能说清慢版本枚举了哪些候选。

\textbf{暴力过程。} 拿 [1, -1, 5, -2, 3]、k=3 手算，到前缀和 3 时需要历史前缀 0；到前缀和 4 时需要历史前缀 1，最早的 1 能给出更长区间。

\textbf{重复工作。} 题面模型：最长区间要求在同一前缀差关系下尽量拉大左右端距离。暴力版的慢点要落到本题：慢版本会为“最长区间要求在同一前缀差关系下尽量拉大左右端距离”反复寻找历史状态；而样例规律已经指出“规律是最长区间要保存每个前缀和的最早下标，而不是最新下标”。这些补数、前缀状态、规范键、半边和或访问过的中间状态可以用一个 key 归类，不必每次重新扫描或重复比较。后面的优化只消掉这类重复工作，不能改变答案集合。

\textbf{优化条件。} 本题的关键观察是：哈希表保存每个前缀和最早出现下标，当前前缀查询 prefix 减 k。这句话必须能翻译成可维护状态；如果题目条件改变后观察不再成立，就不能继续套原来的写法。

\textbf{相邻状态。} 规律是最长区间要保存每个前缀和的最早下标，而不是最新下标。把这个特例继续拆开看：每个合法答案都要还原成当前状态和某个历史状态的配对；哈希表的 key 负责定位历史状态，value 负责表达次数、最早位置或存在性。先查还是先写，必须由是否允许当前前缀和自己配对来决定。

\textbf{状态复用。} 把回头扫描、两两比较或重复状态检测改成查表、分桶或访问记录；key 是当前题目需要的补数、前缀状态、规范表示、半边和或中间状态，value 根据目标选择次数、下标、列表或存在性。本题落点是：规律是最长区间要保存每个前缀和的最早下标，而不是最新下标。手算时只改被新输入影响的那一小部分，而不是回头重算完整候选。

\textbf{指针和字段。} 状态字段要能说清：当前状态或规范键、需要查询的历史 key、哈希表 value 语义、答案累计值或分桶结果；空前缀、空字符串、重复状态或初始状态要单独说明。手算时按这个协议跟踪：拿 [1, -1, 5, -2, 3]、k=3 手算，到前缀和 3 时需要历史前缀 0；到前缀和 4 时需要历史前缀 1，最早的 1 能给出更长区间。然后按协议复查：手算时列出当前输入、当前状态或规范键、需要查询的历史 key、查到的 value、更新后的哈希表。这样能看出先查后插、先分桶后输出，还是先检测重复状态。

\textbf{代码落点。} 先查询还是先插入要由题意决定；key 构造必须无歧义；哈希表 value 的默认插入副作用要可控。关键代码行必须能逐句翻译成“旧状态怎样变成新状态”，否则就只是模板。

\textbf{正确性。} 证明从状态等价开始：任意合法答案都对应当前状态和某个历史 key、规范键或访问状态，查到的 key 也能反推出合法答案或合法分组。

\textbf{边界实验。} 先用从下标零开始、负数、前缀重复时不能覆盖最早位置、答案不存在做反例检查。实验时把从下标零开始、负数、前缀重复时不能覆盖最早位置、答案不存在加入对拍集合，用平方暴力和优化版比较。失败时先看初始历史状态、查询更新顺序和哈希表 value 类型。

\textbf{复杂度变化。} 暴力两两比较、枚举区间、扫描历史或重复迭代常见为平方级或更多；把历史状态、规范键或访问状态放进哈希表后，每步只做常数次查询和更新，通常是线性时间、线性空间。

\textbf{迁移追问。} 如果题目改成“计数版本保存频次，最长版本保存最早位置，这是核心差别”，先判断原状态是否仍然覆盖新答案集合，再决定是微调边界、改 value 语义，还是换算法。

\noindent\textbf{LeetCode 454 四数相加 II。}

\textbf{题目说明。} LeetCode 454 四数相加 II 的题面入口要先还原成具体任务：四个数组的和为零可以拆成两个二数组和互为相反数。

\textbf{样例入口。} 拿 A=[1, 2]、B=[-2, -1]、C=[-1, 2]、D=[0, 2] 手算，先统计 A+B 的和，后面枚举 C+D 时只查相反数。

\textbf{暴力基线。} 慢版本按题意两两比较、枚举区间、扫描历史或持续迭代并查重。把检查条件写出来后，通常能看到当前状态只缺一个历史状态、规范键或已经访问过的中间状态；这正是从平方枚举或重复迭代走向哈希表的入口。放到本题就是：先按“四个数组的和为零可以拆成两个二数组和互为相反数”完整试慢版本；样例里已经能看到“规律是四重枚举拆成两个二重枚举，用哈希表保存一半的和频次”，所以优化前必须先能说清慢版本枚举了哪些候选。

\textbf{暴力过程。} 拿 A=[1, 2]、B=[-2, -1]、C=[-1, 2]、D=[0, 2] 手算，先统计 A+B 的和，后面枚举 C+D 时只查相反数。

\textbf{重复工作。} 题面模型：四个数组的和为零可以拆成两个二数组和互为相反数。暴力版的慢点要落到本题：慢版本会为“四个数组的和为零可以拆成两个二数组和互为相反数”反复寻找历史状态；而样例规律已经指出“规律是四重枚举拆成两个二重枚举，用哈希表保存一半的和频次”。这些补数、前缀状态、规范键、半边和或访问过的中间状态可以用一个 key 归类，不必每次重新扫描或重复比较。后面的优化只消掉这类重复工作，不能改变答案集合。

\textbf{优化条件。} 本题的关键观察是：先统计 A 和 B 的所有两数和频次，再枚举 C 和 D 查询相反数。这句话必须能翻译成可维护状态；如果题目条件改变后观察不再成立，就不能继续套原来的写法。

\textbf{相邻状态。} 规律是四重枚举拆成两个二重枚举，用哈希表保存一半的和频次。把这个特例继续拆开看：每个合法答案都要还原成当前状态和某个历史状态的配对；哈希表的 key 负责定位历史状态，value 负责表达次数、最早位置或存在性。先查还是先写，必须由是否允许当前前缀和自己配对来决定。

\textbf{状态复用。} 把回头扫描、两两比较或重复状态检测改成查表、分桶或访问记录；key 是当前题目需要的补数、前缀状态、规范表示、半边和或中间状态，value 根据目标选择次数、下标、列表或存在性。本题落点是：规律是四重枚举拆成两个二重枚举，用哈希表保存一半的和频次。手算时只改被新输入影响的那一小部分，而不是回头重算完整候选。

\textbf{指针和字段。} 状态字段要能说清：当前状态或规范键、需要查询的历史 key、哈希表 value 语义、答案累计值或分桶结果；空前缀、空字符串、重复状态或初始状态要单独说明。手算时按这个协议跟踪：拿 A=[1, 2]、B=[-2, -1]、C=[-1, 2]、D=[0, 2] 手算，先统计 A+B 的和，后面枚举 C+D 时只查相反数。然后按协议复查：手算时列出当前输入、当前状态或规范键、需要查询的历史 key、查到的 value、更新后的哈希表。这样能看出先查后插、先分桶后输出，还是先检测重复状态。

\textbf{代码落点。} 先查询还是先插入要由题意决定；key 构造必须无歧义；哈希表 value 的默认插入副作用要可控。关键代码行必须能逐句翻译成“旧状态怎样变成新状态”，否则就只是模板。

\textbf{正确性。} 证明从状态等价开始：任意合法答案都对应当前状态和某个历史 key、规范键或访问状态，查到的 key 也能反推出合法答案或合法分组。

\textbf{边界实验。} 先用重复值、负数、答案数量很大、两数和范围做反例检查。实验时把重复值、负数、答案数量很大、两数和范围加入对拍集合，用平方暴力和优化版比较。失败时先看初始历史状态、查询更新顺序和哈希表 value 类型。

\textbf{复杂度变化。} 暴力两两比较、枚举区间、扫描历史或重复迭代常见为平方级或更多；把历史状态、规范键或访问状态放进哈希表后，每步只做常数次查询和更新，通常是线性时间、线性空间。

\textbf{迁移追问。} 如果题目改成“这是 meet-in-the-middle 思想，比四层枚举少两个维度”，先判断原状态是否仍然覆盖新答案集合，再决定是微调边界、改 value 语义，还是换算法。

\noindent\textbf{LeetCode 523 连续的子数组和。}

\textbf{题目说明。} LeetCode 523 连续的子数组和 的题面入口要先还原成具体任务：长度至少二的子数组和为 k 的倍数，等价于两个前缀和同余。

\textbf{样例入口。} 拿 [23, 2, 4, 6, 7]、k=6 手算，前缀余数 5 在下标 0 和 2 重复，中间 [2, 4] 的和能被 6 整除。

\textbf{暴力基线。} 慢版本按题意两两比较、枚举区间、扫描历史或持续迭代并查重。把检查条件写出来后，通常能看到当前状态只缺一个历史状态、规范键或已经访问过的中间状态；这正是从平方枚举或重复迭代走向哈希表的入口。放到本题就是：先按“长度至少二的子数组和为 k 的倍数，等价于两个前缀和同余”完整试慢版本；样例里已经能看到“规律是两个前缀余数相同表示区间和为 k 的倍数，且区间长度至少为 2”，所以优化前必须先能说清慢版本枚举了哪些候选。

\textbf{暴力过程。} 拿 [23, 2, 4, 6, 7]、k=6 手算，前缀余数 5 在下标 0 和 2 重复，中间 [2, 4] 的和能被 6 整除。

\textbf{重复工作。} 题面模型：长度至少二的子数组和为 k 的倍数，等价于两个前缀和同余。暴力版的慢点要落到本题：慢版本会为“长度至少二的子数组和为 k 的倍数，等价于两个前缀和同余”反复寻找历史状态；而样例规律已经指出“规律是两个前缀余数相同表示区间和为 k 的倍数，且区间长度至少为 2”。这些补数、前缀状态、规范键、半边和或访问过的中间状态可以用一个 key 归类，不必每次重新扫描或重复比较。后面的优化只消掉这类重复工作，不能改变答案集合。

\textbf{优化条件。} 本题的关键观察是：哈希表保存某个余数最早出现下标，当前余数若重复且距离至少二则成立。这句话必须能翻译成可维护状态；如果题目条件改变后观察不再成立，就不能继续套原来的写法。

\textbf{相邻状态。} 规律是两个前缀余数相同表示区间和为 k 的倍数，且区间长度至少为 2。把这个特例继续拆开看：每个合法答案都要还原成当前状态和某个历史状态的配对；哈希表的 key 负责定位历史状态，value 负责表达次数、最早位置或存在性。先查还是先写，必须由是否允许当前前缀和自己配对来决定。

\textbf{状态复用。} 把回头扫描、两两比较或重复状态检测改成查表、分桶或访问记录；key 是当前题目需要的补数、前缀状态、规范表示、半边和或中间状态，value 根据目标选择次数、下标、列表或存在性。本题落点是：规律是两个前缀余数相同表示区间和为 k 的倍数，且区间长度至少为 2。手算时只改被新输入影响的那一小部分，而不是回头重算完整候选。

\textbf{指针和字段。} 状态字段要能说清：当前状态或规范键、需要查询的历史 key、哈希表 value 语义、答案累计值或分桶结果；空前缀、空字符串、重复状态或初始状态要单独说明。手算时按这个协议跟踪：拿 [23, 2, 4, 6, 7]、k=6 手算，前缀余数 5 在下标 0 和 2 重复，中间 [2, 4] 的和能被 6 整除。然后按协议复查：手算时列出当前输入、当前状态或规范键、需要查询的历史 key、查到的 value、更新后的哈希表。这样能看出先查后插、先分桶后输出，还是先检测重复状态。

\textbf{代码落点。} 先查询还是先插入要由题意决定；key 构造必须无歧义；哈希表 value 的默认插入副作用要可控。关键代码行必须能逐句翻译成“旧状态怎样变成新状态”，否则就只是模板。

\textbf{正确性。} 证明从状态等价开始：任意合法答案都对应当前状态和某个历史 key、规范键或访问状态，查到的 key 也能反推出合法答案或合法分组。

\textbf{边界实验。} 先用k 为零、负数余数、长度限制、从开头开始的区间做反例检查。实验时把k 为零、负数余数、长度限制、从开头开始的区间加入对拍集合，用平方暴力和优化版比较。失败时先看初始历史状态、查询更新顺序和哈希表 value 类型。

\textbf{复杂度变化。} 暴力两两比较、枚举区间、扫描历史或重复迭代常见为平方级或更多；把历史状态、规范键或访问状态放进哈希表后，每步只做常数次查询和更新，通常是线性时间、线性空间。

\textbf{迁移追问。} 如果题目改成“保存最早下标是为了满足长度条件并保留最长距离”，先判断原状态是否仍然覆盖新答案集合，再决定是微调边界、改 value 语义，还是换算法。

\noindent\textbf{LeetCode 525 连续数组。}

\textbf{题目说明。} LeetCode 525 连续数组 的题面入口要先还原成具体任务：把零看成负一后，零和一数量相等等价于前缀和相同。

\textbf{样例入口。} 拿 [0, 1, 0] 手算，把 0 当 -1，前缀差值序列 -1、0、-1；差值再次出现说明中间 0 和 1 数量相同。

\textbf{暴力基线。} 慢版本按题意两两比较、枚举区间、扫描历史或持续迭代并查重。把检查条件写出来后，通常能看到当前状态只缺一个历史状态、规范键或已经访问过的中间状态；这正是从平方枚举或重复迭代走向哈希表的入口。放到本题就是：先按“把零看成负一后，零和一数量相等等价于前缀和相同”完整试慢版本；样例里已经能看到“规律是保存差值第一次出现的位置，最长区间用最早位置”，所以优化前必须先能说清慢版本枚举了哪些候选。

\textbf{暴力过程。} 拿 [0, 1, 0] 手算，把 0 当 -1，前缀差值序列 -1、0、-1；差值再次出现说明中间 0 和 1 数量相同。

\textbf{重复工作。} 题面模型：把零看成负一后，零和一数量相等等价于前缀和相同。暴力版的慢点要落到本题：慢版本会为“把零看成负一后，零和一数量相等等价于前缀和相同”反复寻找历史状态；而样例规律已经指出“规律是保存差值第一次出现的位置，最长区间用最早位置”。这些补数、前缀状态、规范键、半边和或访问过的中间状态可以用一个 key 归类，不必每次重新扫描或重复比较。后面的优化只消掉这类重复工作，不能改变答案集合。

\textbf{优化条件。} 本题的关键观察是：哈希表保存每个前缀和最早下标，重复出现时更新最长长度。这句话必须能翻译成可维护状态；如果题目条件改变后观察不再成立，就不能继续套原来的写法。

\textbf{相邻状态。} 规律是保存差值第一次出现的位置，最长区间用最早位置。把这个特例继续拆开看：每个合法答案都要还原成当前状态和某个历史状态的配对；哈希表的 key 负责定位历史状态，value 负责表达次数、最早位置或存在性。先查还是先写，必须由是否允许当前前缀和自己配对来决定。

\textbf{状态复用。} 把回头扫描、两两比较或重复状态检测改成查表、分桶或访问记录；key 是当前题目需要的补数、前缀状态、规范表示、半边和或中间状态，value 根据目标选择次数、下标、列表或存在性。本题落点是：规律是保存差值第一次出现的位置，最长区间用最早位置。手算时只改被新输入影响的那一小部分，而不是回头重算完整候选。

\textbf{指针和字段。} 状态字段要能说清：当前状态或规范键、需要查询的历史 key、哈希表 value 语义、答案累计值或分桶结果；空前缀、空字符串、重复状态或初始状态要单独说明。手算时按这个协议跟踪：拿 [0, 1, 0] 手算，把 0 当 -1，前缀差值序列 -1、0、-1；差值再次出现说明中间 0 和 1 数量相同。然后按协议复查：手算时列出当前输入、当前状态或规范键、需要查询的历史 key、查到的 value、更新后的哈希表。这样能看出先查后插、先分桶后输出，还是先检测重复状态。

\textbf{代码落点。} 先查询还是先插入要由题意决定；key 构造必须无歧义；哈希表 value 的默认插入副作用要可控。关键代码行必须能逐句翻译成“旧状态怎样变成新状态”，否则就只是模板。

\textbf{正确性。} 证明从状态等价开始：任意合法答案都对应当前状态和某个历史 key、规范键或访问状态，查到的 key 也能反推出合法答案或合法分组。

\textbf{边界实验。} 先用全零、全一、从下标零开始、不要覆盖最早位置做反例检查。实验时把全零、全一、从下标零开始、不要覆盖最早位置加入对拍集合，用平方暴力和优化版比较。失败时先看初始历史状态、查询更新顺序和哈希表 value 类型。

\textbf{复杂度变化。} 暴力两两比较、枚举区间、扫描历史或重复迭代常见为平方级或更多；把历史状态、规范键或访问状态放进哈希表后，每步只做常数次查询和更新，通常是线性时间、线性空间。

\textbf{迁移追问。} 如果题目改成“这是前缀和最早位置语义，不是频次语义”，先判断原状态是否仍然覆盖新答案集合，再决定是微调边界、改 value 语义，还是换算法。

\noindent\textbf{LeetCode 560 和为 K 的子数组。}

\textbf{题目说明。} LeetCode 560 和为 K 的子数组给定整数数组 nums 和整数 k，统计和等于 k 的连续子数组个数；数组中可以有负数。

\textbf{样例入口。} 样例 nums=[1, 1, 1]，k=2，输出 2；两个合法连续子数组分别是前两个 1 和后两个 1。

\textbf{暴力基线。} 暴力固定 left，向右累计 sum；每当 sum==k 就计数。

\textbf{暴力过程。} 以 nums=[1, 1, 1]、k=2 为例，从 left=0 可找到 [1, 1]，从 left=1 又找到后两个 1。

\textbf{重复工作。} 题面模型：当前前缀和减去目标值，就是需要匹配的历史前缀。暴力版的慢点要落到本题：浪费在每个右端都回头枚举所有 left。区间和等于 k 可以改写成 prefix[right]-prefix[left]=k。后面的优化只消掉这类重复工作，不能改变答案集合。

\textbf{优化条件。} 本题的关键观察是：哈希表保存前缀和出现次数，先查询再更新。这句话必须能翻译成可维护状态；如果题目条件改变后观察不再成立，就不能继续套原来的写法。

\textbf{相邻状态。} 规律是哈希表保存前缀和频次，且初始要放 0:1。把这个特例继续拆开看：每个合法答案都要还原成当前状态和某个历史状态的配对；哈希表的 key 负责定位历史状态，value 负责表达次数、最早位置或存在性。先查还是先写，必须由是否允许当前前缀和自己配对来决定。

\textbf{状态复用。} 扫描前缀和 prefix。哈希表 count[old\_prefix] 保存历史前缀出现次数；当前位置需要查 prefix-k 出现过几次。手算时只改被新输入影响的那一小部分，而不是回头重算完整候选。

\textbf{指针和字段。} 状态字段要能说清：prefix 是当前前缀和，need=prefix-k 是所需历史前缀，count 保存历史前缀频次，answer 累加匹配次数。手算时按这个协议跟踪：[1, 1, 1] 中扫描到第二个元素时 prefix=2，需要历史 0，count[0]=1，所以找到第一个长度 2 子数组。

\textbf{代码落点。} 关键是先 answer += count[prefix-k]，再 count[prefix]++。初始化 count[0]=1，表示从下标 0 开始的区间。关键代码行必须能逐句翻译成“旧状态怎样变成新状态”，否则就只是模板。

\textbf{正确性。} 任意和为 k 的子数组都对应一对前缀差 k；扫描到右端时，所有合法左端前缀都已在表里。

\textbf{边界实验。} 先用负数、目标为零、从下标零开始的区间、重复前缀做反例检查。实验时把负数、目标为零、从下标零开始的区间、重复前缀加入对拍集合，用平方暴力和优化版比较。失败时先看初始历史状态、查询更新顺序和哈希表 value 类型。

\textbf{复杂度变化。} 暴力 O(\ensuremath{n^{2}})，前缀哈希 O(n) 均摊时间、O(n) 空间。

\textbf{迁移追问。} 如果题目改成“若要求最长区间，哈希表应保存最早位置而不是次数”，先判断原状态是否仍然覆盖新答案集合，再决定是微调边界、改 value 语义，还是换算法。

\noindent\textbf{LeetCode 974 和可被 K 整除的子数组。}

\textbf{题目说明。} LeetCode 974 和可被 K 整除的子数组 的题面入口要先还原成具体任务：区间和能被 K 整除等价于两个前缀和余数相同。

\textbf{样例入口。} 拿 [4, 5, 0, -2, -3, 1]、k=5 手算，前缀余数重复时，中间区间和能被 5 整除；负数前缀要先把余数归一化。

\textbf{暴力基线。} 慢版本按题意两两比较、枚举区间、扫描历史或持续迭代并查重。把检查条件写出来后，通常能看到当前状态只缺一个历史状态、规范键或已经访问过的中间状态；这正是从平方枚举或重复迭代走向哈希表的入口。放到本题就是：先按“区间和能被 K 整除等价于两个前缀和余数相同”完整试慢版本；样例里已经能看到“规律是哈希表保存余数频次，每到一个余数就增加已有次数”，所以优化前必须先能说清慢版本枚举了哪些候选。

\textbf{暴力过程。} 拿 [4, 5, 0, -2, -3, 1]、k=5 手算，前缀余数重复时，中间区间和能被 5 整除；负数前缀要先把余数归一化。

\textbf{重复工作。} 题面模型：区间和能被 K 整除等价于两个前缀和余数相同。暴力版的慢点要落到本题：慢版本会为“区间和能被 K 整除等价于两个前缀和余数相同”反复寻找历史状态；而样例规律已经指出“规律是哈希表保存余数频次，每到一个余数就增加已有次数”。这些补数、前缀状态、规范键、半边和或访问过的中间状态可以用一个 key 归类，不必每次重新扫描或重复比较。后面的优化只消掉这类重复工作，不能改变答案集合。

\textbf{优化条件。} 本题的关键观察是：统计每个余数出现次数，当前余数能和所有历史同余前缀组成答案。这句话必须能翻译成可维护状态；如果题目条件改变后观察不再成立，就不能继续套原来的写法。

\textbf{相邻状态。} 规律是哈希表保存余数频次，每到一个余数就增加已有次数。把这个特例继续拆开看：每个合法答案都要还原成当前状态和某个历史状态的配对；哈希表的 key 负责定位历史状态，value 负责表达次数、最早位置或存在性。先查还是先写，必须由是否允许当前前缀和自己配对来决定。

\textbf{状态复用。} 把回头扫描、两两比较或重复状态检测改成查表、分桶或访问记录；key 是当前题目需要的补数、前缀状态、规范表示、半边和或中间状态，value 根据目标选择次数、下标、列表或存在性。本题落点是：规律是哈希表保存余数频次，每到一个余数就增加已有次数。手算时只改被新输入影响的那一小部分，而不是回头重算完整候选。

\textbf{指针和字段。} 状态字段要能说清：当前状态或规范键、需要查询的历史 key、哈希表 value 语义、答案累计值或分桶结果；空前缀、空字符串、重复状态或初始状态要单独说明。手算时按这个协议跟踪：拿 [4, 5, 0, -2, -3, 1]、k=5 手算，前缀余数重复时，中间区间和能被 5 整除；负数前缀要先把余数归一化。然后按协议复查：手算时列出当前输入、当前状态或规范键、需要查询的历史 key、查到的 value、更新后的哈希表。这样能看出先查后插、先分桶后输出，还是先检测重复状态。

\textbf{代码落点。} 先查询还是先插入要由题意决定；key 构造必须无歧义；哈希表 value 的默认插入副作用要可控。关键代码行必须能逐句翻译成“旧状态怎样变成新状态”，否则就只是模板。

\textbf{正确性。} 证明从状态等价开始：任意合法答案都对应当前状态和某个历史 key、规范键或访问状态，查到的 key 也能反推出合法答案或合法分组。

\textbf{边界实验。} 先用负数余数归一化、K 为一、从开头开始、重复余数很多做反例检查。实验时把负数余数归一化、K 为一、从开头开始、重复余数很多加入对拍集合，用平方暴力和优化版比较。失败时先看初始历史状态、查询更新顺序和哈希表 value 类型。

\textbf{复杂度变化。} 暴力两两比较、枚举区间、扫描历史或重复迭代常见为平方级或更多；把历史状态、规范键或访问状态放进哈希表后，每步只做常数次查询和更新，通常是线性时间、线性空间。

\textbf{迁移追问。} 如果题目改成“如果问最长长度，保存最早位置而不是次数”，先判断原状态是否仍然覆盖新答案集合，再决定是微调边界、改 value 语义，还是换算法。

\noindent\textbf{LeetCode 1074 元素和为目标值的子矩阵数量。}

\textbf{题目说明。} LeetCode 1074 元素和为目标值的子矩阵数量 的题面入口要先还原成具体任务：二维矩阵子区域和可以固定上下边界后转成一维子数组和。

\textbf{样例入口。} 拿 3x3 矩阵和 target=0 手算，先固定上下两行，把每列压成列和，问题变成一维子数组和为 0 的计数。

\textbf{暴力基线。} 慢版本按题意两两比较、枚举区间、扫描历史或持续迭代并查重。把检查条件写出来后，通常能看到当前状态只缺一个历史状态、规范键或已经访问过的中间状态；这正是从平方枚举或重复迭代走向哈希表的入口。放到本题就是：先按“二维矩阵子区域和可以固定上下边界后转成一维子数组和”完整试慢版本；样例里已经能看到“规律是枚举行边界，列方向用前缀和加哈希计数，二维问题降成多次一维问题”，所以优化前必须先能说清慢版本枚举了哪些候选。

\textbf{暴力过程。} 拿 3x3 矩阵和 target=0 手算，先固定上下两行，把每列压成列和，问题变成一维子数组和为 0 的计数。

\textbf{重复工作。} 题面模型：二维矩阵子区域和可以固定上下边界后转成一维子数组和。暴力版的慢点要落到本题：慢版本会为“二维矩阵子区域和可以固定上下边界后转成一维子数组和”反复寻找历史状态；而样例规律已经指出“规律是枚举行边界，列方向用前缀和加哈希计数，二维问题降成多次一维问题”。这些补数、前缀状态、规范键、半边和或访问过的中间状态可以用一个 key 归类，不必每次重新扫描或重复比较。后面的优化只消掉这类重复工作，不能改变答案集合。

\textbf{优化条件。} 本题的关键观察是：枚举行边界，压缩每列区间和，再用前缀和哈希统计目标和子数组。这句话必须能翻译成可维护状态；如果题目条件改变后观察不再成立，就不能继续套原来的写法。

\textbf{相邻状态。} 规律是枚举行边界，列方向用前缀和加哈希计数，二维问题降成多次一维问题。把这个特例继续拆开看：每个合法答案都要还原成当前状态和某个历史状态的配对；哈希表的 key 负责定位历史状态，value 负责表达次数、最早位置或存在性。先查还是先写，必须由是否允许当前前缀和自己配对来决定。

\textbf{状态复用。} 把回头扫描、两两比较或重复状态检测改成查表、分桶或访问记录；key 是当前题目需要的补数、前缀状态、规范表示、半边和或中间状态，value 根据目标选择次数、下标、列表或存在性。本题落点是：规律是枚举行边界，列方向用前缀和加哈希计数，二维问题降成多次一维问题。手算时只改被新输入影响的那一小部分，而不是回头重算完整候选。

\textbf{指针和字段。} 状态字段要能说清：当前状态或规范键、需要查询的历史 key、哈希表 value 语义、答案累计值或分桶结果；空前缀、空字符串、重复状态或初始状态要单独说明。手算时按这个协议跟踪：拿 3x3 矩阵和 target=0 手算，先固定上下两行，把每列压成列和，问题变成一维子数组和为 0 的计数。然后按协议复查：手算时列出当前输入、当前状态或规范键、需要查询的历史 key、查到的 value、更新后的哈希表。这样能看出先查后插、先分桶后输出，还是先检测重复状态。

\textbf{代码落点。} 先查询还是先插入要由题意决定；key 构造必须无歧义；哈希表 value 的默认插入副作用要可控。关键代码行必须能逐句翻译成“旧状态怎样变成新状态”，否则就只是模板。

\textbf{正确性。} 证明从状态等价开始：任意合法答案都对应当前状态和某个历史 key、规范键或访问状态，查到的 key 也能反推出合法答案或合法分组。

\textbf{边界实验。} 先用单行、单列、负数、目标为零、矩阵较大做反例检查。实验时把单行、单列、负数、目标为零、矩阵较大加入对拍集合，用平方暴力和优化版比较。失败时先看初始历史状态、查询更新顺序和哈希表 value 类型。

\textbf{复杂度变化。} 暴力两两比较、枚举区间、扫描历史或重复迭代常见为平方级或更多；把历史状态、规范键或访问状态放进哈希表后，每步只做常数次查询和更新，通常是线性时间、线性空间。

\textbf{迁移追问。} 如果题目改成“这是二维前缀思想和一维哈希计数的组合”，先判断原状态是否仍然覆盖新答案集合，再决定是微调边界、改 value 语义，还是换算法。

\noindent\textbf{LeetCode 1248 统计优美子数组。}

\textbf{题目说明。} LeetCode 1248 统计优美子数组 的题面入口要先还原成具体任务：恰好 K 个奇数可转成前缀奇数计数差为 K。

\textbf{样例入口。} 拿 [1, 1, 2, 1, 1]、k=3 手算，奇数个数前缀从 0 到 3 时，需要历史奇数个数 0；每个匹配历史前缀对应一个优美子数组。

\textbf{暴力基线。} 慢版本按题意两两比较、枚举区间、扫描历史或持续迭代并查重。把检查条件写出来后，通常能看到当前状态只缺一个历史状态、规范键或已经访问过的中间状态；这正是从平方枚举或重复迭代走向哈希表的入口。放到本题就是：先按“恰好 K 个奇数可转成前缀奇数计数差为 K”完整试慢版本；样例里已经能看到“规律是把奇数看成 1、偶数看成 0，转成前缀和计数”，所以优化前必须先能说清慢版本枚举了哪些候选。

\textbf{暴力过程。} 拿 [1, 1, 2, 1, 1]、k=3 手算，奇数个数前缀从 0 到 3 时，需要历史奇数个数 0；每个匹配历史前缀对应一个优美子数组。

\textbf{重复工作。} 题面模型：恰好 K 个奇数可转成前缀奇数计数差为 K。暴力版的慢点要落到本题：慢版本会为“恰好 K 个奇数可转成前缀奇数计数差为 K”反复寻找历史状态；而样例规律已经指出“规律是把奇数看成 1、偶数看成 0，转成前缀和计数”。这些补数、前缀状态、规范键、半边和或访问过的中间状态可以用一个 key 归类，不必每次重新扫描或重复比较。后面的优化只消掉这类重复工作，不能改变答案集合。

\textbf{优化条件。} 本题的关键观察是：哈希表保存某个奇数计数出现次数，当前计数查询 count - K。这句话必须能翻译成可维护状态；如果题目条件改变后观察不再成立，就不能继续套原来的写法。

\textbf{相邻状态。} 规律是把奇数看成 1、偶数看成 0，转成前缀和计数。把这个特例继续拆开看：每个合法答案都要还原成当前状态和某个历史状态的配对；哈希表的 key 负责定位历史状态，value 负责表达次数、最早位置或存在性。先查还是先写，必须由是否允许当前前缀和自己配对来决定。

\textbf{状态复用。} 把回头扫描、两两比较或重复状态检测改成查表、分桶或访问记录；key 是当前题目需要的补数、前缀状态、规范表示、半边和或中间状态，value 根据目标选择次数、下标、列表或存在性。本题落点是：规律是把奇数看成 1、偶数看成 0，转成前缀和计数。手算时只改被新输入影响的那一小部分，而不是回头重算完整候选。

\textbf{指针和字段。} 状态字段要能说清：当前状态或规范键、需要查询的历史 key、哈希表 value 语义、答案累计值或分桶结果；空前缀、空字符串、重复状态或初始状态要单独说明。手算时按这个协议跟踪：拿 [1, 1, 2, 1, 1]、k=3 手算，奇数个数前缀从 0 到 3 时，需要历史奇数个数 0；每个匹配历史前缀对应一个优美子数组。然后按协议复查：手算时列出当前输入、当前状态或规范键、需要查询的历史 key、查到的 value、更新后的哈希表。这样能看出先查后插、先分桶后输出，还是先检测重复状态。

\textbf{代码落点。} 先查询还是先插入要由题意决定；key 构造必须无歧义；哈希表 value 的默认插入副作用要可控。关键代码行必须能逐句翻译成“旧状态怎样变成新状态”，否则就只是模板。

\textbf{正确性。} 证明从状态等价开始：任意合法答案都对应当前状态和某个历史 key、规范键或访问状态，查到的 key 也能反推出合法答案或合法分组。

\textbf{边界实验。} 先用K 为零、全偶数、全奇数、从开头开始的区间做反例检查。实验时把K 为零、全偶数、全奇数、从开头开始的区间加入对拍集合，用平方暴力和优化版比较。失败时先看初始历史状态、查询更新顺序和哈希表 value 类型。

\textbf{复杂度变化。} 暴力两两比较、枚举区间、扫描历史或重复迭代常见为平方级或更多；把历史状态、规范键或访问状态放进哈希表后，每步只做常数次查询和更新，通常是线性时间、线性空间。

\textbf{迁移追问。} 如果题目改成“也可用至多 K 减至多 K-1 的滑动窗口”，先判断原状态是否仍然覆盖新答案集合，再决定是微调边界、改 value 语义，还是换算法。

\subsection{实验复盘}

追问通常会问先查询还是先更新、哈希表保存次数还是最早下标、为什么负数不影响前缀和。请准备说明从任意合法区间如何反推到两个前缀状态。

复盘本节时先从 LeetCode 1 两数之和、LeetCode 49 字母异位词分组、LeetCode 128 最长连续序列 开始：每道题都先手写一个能覆盖全部候选的暴力版，再指出优化结构具体保存了哪一段历史信息，最后用相邻案例比较为什么不能互换解法。
